\section{Online Accessibility and Academic Accommodations}\label{online-accessibility-and-academic-accommodations}

\stanceinfo{\textbf{CELC 2024} [2024-01-05]}{2024-01-05}

\officialstance{The Canadian Federation of Engineering Students believes that educational resources [recordings, lecture notes etc.] should remain accessible online at the same level of quality as in-person sessions. The CFES further believes that accommodations need to be in place for assessments throughout the academic term for extenuating circumstances, illness, and other personal matters, providing catered support to all engineering students.}

\stanceheading{The Students' Position}
\begin{itemize}
 \item All engineering students should have equal opportunity to pursue and excel in their education, regardless of individual abilities and circumstances.
 \item COVID-19 presented the opportunity to have permanent online resources that have been favored by engineering students to improve their academic performance.
 \item A universal design approach for educational resources and examinations does not accommodate for diverse needs and learning challenges.
 \item Traditional assessments do not reflect understanding of content in all situations due to the high stress environment and necessity to perform on the spot.
 \item Accommodations should be provided on a need basis to cater and support engineering students in their education and to encourage their success.
\end{itemize}

\stanceheading{Recommendations to Partners, Stakeholders, and Other Entities}
\begin{itemize}
 \item The CFES calls on the Canadian Engineering Accreditation Board (CEAB) to encourage institutions in the process of updating their institution policies and being more inclusive to hybrid education and online permanent materials.
 \item The CFES calls on the Canadian Engineering Accreditation Board to assess the accommodations and accessibility available in an engineering program when undergoing an accreditation visit.
 \item The CFES calls on the Engineering Deans Canada to encourage institutions to create a standard hybrid model for course delivery based on student feedback on the program.
 \item The CFES calls on the Engineering Deans Canada to ensure accommodations are granted to those with diverse circumstances or disadvantages in order to promote equity in engineering.
\end{itemize}

\stanceheading{What the CFES is doing}
\begin{itemize}
 \item The 2023 National Survey is collecting information on academic resources and availability.
 \item Literature reviews on education in a post-pandemic world that found positives from online learning.
 \item The CFES continues to meet with stakeholders surrounding the flexibility of education and quality of education.
\end{itemize}

\stanceheading{What the CFES plans to do}
\begin{itemize}
 \item The CFES will work with stakeholders and partners to find an even ground to begin policy work for online accommodations and academic accessibility in engineering education.
 \item The CFES will create resources to support the membership to advocate for online accessibility and academic accommodations through letters of support.
 \item The CFES will gather more information and testimonials on current student policy and experiences surrounding online accessibility and academic accommodations.
\end{itemize}

\newpage
